%% =========================================================================
\section{Methods}\label{sec:methods}

The demolition indicators D1--D6 (Section~\ref{sec:indicators}) are
applied to the cleaned extract of Section~\ref{sec:data}. This section
describes the remaining analysis choices: the treatment of discontinued
use codes, the floor-area measure, the quantities computed
for each indicator, and the planned comparisons.

\subsection{Treatment of discontinued use codes}\label{sec:discontinued-method}

A fixed set of retired, round-number building-use codes (agricultural
and industrial categories no longer issued) is associated with an
ambiguous status-10 signal: earlier work found that a status-10 event on
these codes sometimes reflects administrative re-registration rather
than physical demolition. Rather than building this into a separate
indicator, it is implemented as an orthogonal, on/off exclusion axis
applicable to any of D1--D6 (Section~\ref{sec:indicators}); applying it
to D1 reproduces the earlier inclusive/filtered contrast directly (D1:
436{,}194 buildings; D1 with the exclusion applied: 302{,}261). Whether
excluding these codes should be treated
as a correction is, on the current extract, an open empirical question
rather than an assumption: the exclusion strips approximately 34\,\% of
every indicator's buildings uniformly, consistent with a base rate
across the stock rather than concentrated artifacts, and among the
status-10 buildings it would remove, 50.7\,\% also carry a formal
demolition case, against a 44.2\,\% case-rate among the buildings it
would keep --- so on this extract the exclusion also removes buildings
with independent demolition evidence. Crossing every indicator with the
axis on and off gives the full $7\times2$ comparison grid used in the
ablation (Section~\ref{sec:comparisons}); whether the exclusion improves
or worsens agreement with ground truth is intended to be settled against
BOSSINF (Section~\ref{sec:comparisons}), not assumed.

\subsection{Floor-area measure}\label{sec:areas}

The primary area outcome is total building area,
\texttt{byg038SamletBygningsareal}. This is the area concept used for the
main demolished-square-metre totals and, where a stock denominator is
available, for demolition rates by area. We use a single primary area
field so that the ablation isolates the demolition indicator: the
building set changes from D1 to D6, while the area definition is held
fixed.

This choice is also a comparability decision. \texttt{byg038} is BBR's
registered total floor area and is the closest single field to the area
concept used in service-life, building-stock, and material-flow
applications. Its limitation is coverage, not interpretation:
Section~\ref{sec:missingness} shows that \texttt{byg038} is missing for a
large and non-random part of the stock, especially small outbuildings.
Therefore a building with missing total area remains in the demolition
count, but contributes no square metres to the area total; the number and
share of buildings with non-missing total area are reported alongside the
area estimates. Negative \texttt{byg038} values are excluded after
candidate-event detection using the same rule for every indicator. Zero
values are retained as recorded zero-area observations and reported in
the coverage diagnostics rather than being treated as positive demolished
area.

Other BBR area fields are retained as diagnostics and sensitivity checks,
but not as co-equal headline outcomes. The footprint
\texttt{byg041BebyggetAreal} is nearly complete for historical buildings,
but it measures ground coverage rather than floor area and is therefore
not a substitute for \texttt{byg038}. The sum of residential and
commercial area (\texttt{byg039}+\texttt{byg040}) is useful for comparison
with BUILD's published analysis~\cite{jensen2025nedrivning}, but it is a
different constructed area definition. These alternatives are reported as
robustness or appendix material rather than used to define the main
indicator comparison. A hybrid or imputed area measure---for example,
filling missing total area from footprint and storeys---could be useful
for estimating a more complete physical area flow, but it would introduce
an additional modelling assumption and is therefore left outside the main
indicator-sensitivity analysis.

\subsection{Computed quantities}\label{sec:metrics}

For each indicator, the following are computed by year: the number of
demolitions; the demolished total area; the average and median demolished
area per building; the coverage of \texttt{byg038}; and the demolition
rate relative to the building stock where a valid denominator can be
constructed. The same filters and aggregation logic are used for all
indicators. Results are computed nationally, and by building-use group,
region, and construction period where counts allow. For comparability of
disaggregated results, use codes are grouped according to the 22-group
mapping used by BUILD~\cite{jensen2025nedrivning}.
\groupdecision{exact headline metric (annual demolished m$^2$; share of
stock m$^2$ per year; count-based rates) --- open question \#1 in the
project notes.}

As defined in Section~\ref{sec:data}, annual series are produced over
the full reporting window (2000--2025), while headline averages,
windowed overlap, and rate comparisons use the main comparison window
(2018--2025). The KMD comparison uses its own external-anchor window
(2011--2019).

\subsection{Planned comparisons}\label{sec:comparisons}

%% From the analysis plan (docs/plan.md, Analyses 2, 7, 8): indicator
%% overlap; BUILD reproduction; overlap with the independent extract.
Three comparisons are planned alongside the scenario grid.
\emph{Indicator agreement:} which buildings (and how much area) each
indicator identifies, pairwise and jointly.
\emph{BUILD:} a reproduction of BUILD's analysis setup---buildings
constructed before 1999, retired 2012--2023, area as residential plus
commercial floor area---against their published totals
(Section~\ref{sec:build-comparison})~\cite{jensen2025nedrivning}.
\emph{Independent demolition data:} overlap with an independently
compiled national extract of reported demolition cases (2000--2020)
previously used in published service-life
research~\cite{andersen2022adaptation, andersen2023lifespan}, treated as
a comparison source rather than ground truth, and with the register of
grant-funded demolitions used in the evaluation of the village-renewal
demolition subsidies~\cite{jensen2019landsbyfornyelse}.
\groupdecision{final scope of the validation comparisons --- which of
these are in the paper versus supplementary material.}
